\section{Kungfu Prototype and Design}
\label{sec:design}

Kungfu is infrastructure for managing real-world work with agents. In this
context, an observer-declared timeline is one surface of a broader runtime fact
ledger. The current implementation establishes the Episode substrate beneath
that surface; it does not yet implement the complete multi-machine observer
projection.

\subsection{Runtime fact ledger}

The runtime fact ledger preserves source-local facts, provenance, causal links,
payload evidence, schema bindings, and verification results. Future sync views
also need accepted source ranges and freshness boundaries. The visible timeline
is a projection over the ledger, not the ledger itself.

This distinction matters. If storage flattens remote facts into anonymous local
records, the product loses the ability to explain what was accepted and why. If
the GUI presents a global-looking order without observer metadata, the user or
agent cannot reproduce the view during recovery.

\subsection{Implemented Episode path}

The Episode prototype separates five layers \cite{kungfu2026episode}:

\begin{enumerate}[leftmargin=*]
  \item An append-only yijinjing manifest journal stores fixed-layout Episode
        lifecycle, frame, reference, and dependency records as authority.
  \item One typed C++ fold derives the canonical current Episode view. Python,
        Node, CLI, and GUI edges do not define independent Episode semantics.
  \item Content-addressed immutable storage resolves payload and schema
        references, with hashes and lengths verified before use.
  \item Structural fsck checks manifest integrity, causal closure, dependencies,
        content evidence, lifecycle state, and projection drift. Recovery and
        repair append evidence instead of silently rewriting verified history.
  \item A SQLite Episode projection accelerates query but remains rebuildable
        from the journal. Absence or drift degrades the projection without
        promoting the database to authority.
\end{enumerate}

This architecture provides the accepted-object and evidence substrate required
by an observer-declared timeline. It also makes partial failure visible at the
Episode boundary: lifecycle, health, and operation capability are separate
dimensions.

\subsection{KFD-4 implementation map}

Table~\ref{tab:kfd4-map} distinguishes implemented evidence from the remaining
observer-view work.

\begin{table}[ht]
\centering
\small
\begin{tabular}{>{\raggedright\arraybackslash}p{0.22\linewidth}>{\raggedright\arraybackslash}p{0.34\linewidth}>{\raggedright\arraybackslash}p{0.34\linewidth}}
\toprule
KFD-4 field & Current substrate & Remaining view work \\
\midrule
Observer & Source and writer-location provenance & First-class observer and stable view identifier \\
Accepted facts & Episode manifests, refs, dependencies, export/import evidence & Multi-source accepted ranges, watermarks, and freshness \\
Projection policy & Deterministic typed fold and rebuildable Episode query projection & Versioned cross-source ordering policy and tie-breaker \\
Causal constraints & In-Episode closure and declared Episode dependencies & Projection-level enforcement over a selected multi-source set \\
Degraded evidence & Structured fsck issues, missing evidence, stale projection, recovery paths & View-level degradation for incomplete ranges and policy knowledge \\
Verification & Semantic oracle, fsck, Trust Report, CI gate & Exported observer declaration and projection-level fsck \\
\bottomrule
\end{tabular}
\caption{Current Kungfu evidence mapped to the KFD-4 observer-perspective contract.}
\label{tab:kfd4-map}
\end{table}

\subsection{Dual-first interfaces}

The same timeline contract should be usable by humans and agents. A GUI can
show the default observer view while exposing metadata on demand. A CLI or API
can return the same accepted facts, projection policy, and degraded evidence
state directly. Agent-facing surfaces are not secondary integrations; they are
first-class control surfaces for participants that may inspect the system more
frequently than humans.

\subsection{Remote work}

Real-world agent work often happens on remote machines. A server-side agent may
write facts into a remote Kungfu home while the local GUI remains the user's
primary control plane. A useful product should support importing or syncing
remote facts back to the local perspective without pretending that the local
observer saw every event live. The imported timeline should declare source
location, accepted range, freshness, and any degraded evidence.

\subsection{Verification gates}

KFD-4 can be used as a release or product gate for perspective-bearing
surfaces. Passing the gate should mean that the surface exposes the observer,
accepted facts, projection policy, causal constraints, and degraded evidence
metadata. It should not mean that every adapter, remote source, or runtime
payload is complete. Stronger claims require their own KFD-2-style trust
assessment.

The Episode qualification report is an example of that evidence discipline. It
records source revision, profile, workload, semantic dimensions, violations,
performance observations, and known gaps. It is evidence for a bounded Episode
profile, not an automatic KFD-4 certificate for every timeline surface.
